\documentclass[11pt,a4paper]{article}
% Drake_preamble.tex - LaTeX Preamble Template
% 作者: 謝慕揚, MD, PhD, FESC
% 
% 使用說明:
% 1. 表格請使用 longtable，不要有垂直分隔線 |
% 2. 表格內換行使用 \newline，不要使用 \\
% 3. 藥物名稱、檢驗、檢查請維持英文
% 4. Reference 格式請使用 \href 加上驗證過的 DOI

% Including essential packages for Chinese typesetting
\usepackage{xeCJK}
\usepackage{fontspec}
% Configuring Chinese font with Noto Serif CJK TC, fallback to SimSun
\IfFontExistsTF{Noto Serif CJK TC}{
  \setCJKmainfont{Noto Serif CJK TC}
  \setCJKsansfont{Noto Serif CJK TC}
  \setCJKmonofont{Noto Serif CJK TC}
}{\setCJKmainfont{SimSun}
  \setCJKsansfont{SimSun}
  \setCJKmonofont{SimSun}
}

% Including other necessary packages
\usepackage{geometry}
\usepackage{titlesec}
\usepackage{enumitem}
\usepackage{xcolor}
\usepackage{colortbl}
\usepackage{tcolorbox}
\usepackage{booktabs}
\usepackage{longtable}
\usepackage{array}
\usepackage{graphicx}
\usepackage{tikz}
\usepackage{fancyhdr}
\usepackage{multicol}
\usepackage{datetime2}
\usepackage{hyperref}
\usepackage{mdframed}
\usepackage{amsmath}
\usepackage{amssymb}
\usepackage{multirow}

\hypersetup{
    colorlinks=true,
    linkcolor=emergency_blue,
    citecolor=emergency_blue,
    urlcolor=emergency_blue,
    pdfborder={0 0 0}
}

% Defining page geometry
\geometry{
  top=2.5cm,
  bottom=2.5cm,
  left=2cm,
  right=2cm,
  headheight=25pt
}

% Adjusting topmargin to compensate for increased headheight
\addtolength{\topmargin}{-10pt}

% Defining colors
\definecolor{emergency_red}{RGB}{186,24,27}
\definecolor{emergency_blue}{RGB}{0,114,188}
\definecolor{emergency_gray}{RGB}{120,120,120}
\definecolor{highlight_yellow}{RGB}{255,250,205}

% Setting title formats
\titleformat{\section}[block]{\Large\bfseries\color{emergency_red}}{\thesection}{10pt}{}
\titleformat{\subsection}[block]{\large\bfseries\color{emergency_blue}}{\thesubsection}{8pt}{}
\titleformat{\subsubsection}[block]{\normalsize\bfseries\color{emergency_gray}}{\thesubsubsection}{6pt}{}

% Defining custom environment for important notes
\newmdenv[
    linecolor=emergency_red,
    backgroundcolor=highlight_yellow,
    linewidth=2pt,
    topline=true,
    bottomline=true,
    leftline=true,
    rightline=true,
    innertopmargin=10pt,
    innerbottommargin=10pt,
    innerleftmargin=10pt,
    innerrightmargin=10pt
]{importantbox}

% Defining note command with tcolorbox
\newcommand{\note}[1]{
    \par
    \noindent
    \begin{tcolorbox}[
        colback=emergency_blue!5!white,
        colframe=emergency_blue,
        title=\textbf{重點提醒},
        fonttitle=\bfseries,
        boxrule=0.5pt,
        arc=3pt,
        left=6pt,
        right=6pt,
        top=6pt,
        bottom=6pt
    ]
    #1
    \end{tcolorbox}
}

% Key findings box
\newcommand{\keyfinding}[1]{
    \par
    \noindent
    \begin{tcolorbox}[
        colback=yellow!10!white,
        colframe=orange!75!black,
        title=\textbf{核心發現摘要},
        fonttitle=\bfseries,
        boxrule=0.5pt,
        arc=3pt
    ]
    #1
    \end{tcolorbox}
}

% Defining custom date format for ISO 8601
\DTMnewdatestyle{iso}{
  \renewcommand{\DTMdisplaydate}[4]{\number##1-\number##2-\number##3}
  \renewcommand{\DTMDisplaydate}[4]{\number##1-\number##2-\number##3}
}
\DTMsetdatestyle{iso}
\newcommand{\mydate}{\DTMtoday}


% Custom header
\pagestyle{fancy}
\fancyhf{}
\fancyhead[L]{\textcolor{emergency_red}{ICM 教學講義}}
\fancyhead[R]{\textcolor{emergency_blue}{Cardiac Output Monitoring}}
\fancyfoot[C]{\textcolor{emergency_gray}{\thepage}}
\renewcommand{\headrulewidth}{1pt}
\renewcommand{\headrule}{\hbox to\headwidth{\color{emergency_red}\leaders\hrule height \headrulewidth\hfill}}

\begin{document}

%============================================================
% Title Page
%============================================================
\begin{center}
{\LARGE\bfseries\textcolor{emergency_red}{Intensive Care Medicine}}\\[0.5cm]
{\Large\textbf{Monitoring Cardiac Output}}\\[0.3cm]
{\large 心輸出量監測：Toolbox Article}\\[0.5cm]
{\normalsize \href{https://doi.org/10.1007/s00134-024-07566-6}{Intensive Care Med (2024) 50:1912-1915. DOI: 10.1007/s00134-024-07566-6}}\\[0.3cm]
{\normalsize 整理：謝慕揚 MD, PhD, FESC \quad 日期：\mydate}
\end{center}

\vspace{0.5cm}

%============================================================
\section{前言}
%============================================================

Cardiac output (CO) 監測是重症患者血流動力學評估的核心。CO 整合了 heart rate、preload、afterload 及左右心室功能，反映心臟的泵血能力，並決定組織氧氣輸送。臨床評估無法準確估計 CO 值或其變化，因此應作為血流動力學評估的觸發點而非替代品。

\keyfinding{
\begin{itemize}
    \item CO 監測適應症：初始輸液復甦及小劑量升壓劑後仍無反應的循環衰竭患者
    \item Echocardiography 提供完整心臟評估，但無法作為唯一連續監測工具
    \item 選擇監測工具需考量侵入性及患者狀況
    \item CO 變化需超過 least significant change 才具臨床意義
\end{itemize}
}

%============================================================
\section{CO 測量方法}
%============================================================

\subsection{Echocardiography}

\textbf{原理}：從 left ventricular outflow tract (LVOT) 直徑及 time-averaged velocity of flow (VTI\textsubscript{LVOT}) 計算 CO。

\begin{longtable}{p{4cm}p{10cm}}
\toprule
\textbf{項目} & \textbf{內容} \\
\midrule
\endfirsthead
\toprule
\textbf{項目} & \textbf{內容} \\
\midrule
\endhead
優點 & 非侵入性\newline 快速設置\newline 準確可靠\newline 完整血流動力學評估（morphology、pressures、volumes、systolic/diastolic function）\newline 可評估 fluid responsiveness \\
\midrule
缺點 & 間歇性測量\newline 需技術熟練\newline 耗時（難以每日重複超過兩次）\newline 影像品質有時無法測量所有指標 \\
\midrule
注意事項 & 對 LVOT 直徑及 VTI\textsubscript{LVOT} 測量誤差敏感\newline 為減少觀察者間變異，常使用 VTI\textsubscript{LVOT}（反映 stroke volume）而非直接測量 CO \\
\bottomrule
\end{longtable}

\subsection{Pulse-Derived CO}

\textbf{原理}：動脈脈波受 stroke volume、aortic elastance 及 vascular tone 影響。專用演算法分析動脈波形以估計 stroke volume。

\subsubsection{Non-Calibrated / Internally Calibrated Devices}

\begin{itemize}
    \item 使用複雜演算法估計 arterial tone
    \item 可分析標準 radial catheter 或非侵入性 finger-cuff 裝置
    \item 可在各種血流動力學條件下估計 CO
    \item \textcolor{emergency_red}{限制}：血管張力急性變化時不可靠
    \item 僅提供 CO 及 stroke volume variation (SVV)
\end{itemize}

\subsubsection{Externally Calibrated Devices（Transpulmonary Thermodilution）}

\begin{itemize}
    \item 更準確估計 CO
    \item 提供完整血流動力學評估：cardiac preload、lung water and permeability、cardiac contractility
    \item \textcolor{emergency_blue}{校正時機}：超過 1 小時未校正或血管張力改變後需重新校正
    \item 需專用動脈導管
\end{itemize}

\subsection{Right-Sided Thermodilution（Pulmonary Artery Catheter, PAC）}

\textbf{原理}：標準化注射冷鹽水，根據溫度變化計算 CO。

\begin{longtable}{p{4cm}p{10cm}}
\toprule
\textbf{類型} & \textbf{特點} \\
\midrule
\endfirsthead
\toprule
\textbf{類型} & \textbf{特點} \\
\midrule
\endhead
傳統 PAC & 間歇性 bolus thermodilution\newline 準確可靠 \\
\midrule
改良型（heating filament） & 半連續自動測量\newline 因 2-3 分鐘平均化延遲，無法偵測急性 CO 變化 \\
\midrule
新一代（pulse wave analysis） & Real-time CO 測量\newline 以 bolus thermodilution 校正\newline 需更多驗證 \\
\bottomrule
\end{longtable}

\textbf{PAC 提供的額外資訊}：
\begin{itemize}
    \item Pulmonary artery pressures
    \item Pulmonary artery occluded pressure (PAOP)（左心房壓替代指標）
    \item Right atrial pressure
    \item Mixed-venous O\textsubscript{2} saturation (SvO\textsubscript{2})
\end{itemize}

%============================================================
\section{各種技術比較}
%============================================================

\begin{longtable}{p{4cm}p{5cm}p{5cm}}
\toprule
\textbf{技術} & \textbf{主要優點} & \textbf{主要缺點} \\
\midrule
\endfirsthead
\toprule
\textbf{技術} & \textbf{主要優點} & \textbf{主要缺點} \\
\midrule
\endhead
Echocardiography & 非侵入性\newline 完整心臟評估 & 間歇性\newline 需技術熟練 \\
\midrule
Bioimpedance & 非侵入性\newline 連續 & 重症患者不可靠 \\
\midrule
Bioreactance & 非侵入性\newline 連續 & 重症患者 CO 測量不準確 \\
\midrule
Esophageal Doppler & 微創\newline 連續 & 需頻繁重新定位\newline 非鎮靜患者不適 \\
\midrule
End-tidal CO\textsubscript{2} & 連續\newline 便宜 & 僅追蹤變化，不估計絕對值\newline 需侵入性控制呼吸 \\
\midrule
Non-calibrated PWA & 連續\newline 可評估 fluid responsiveness & CO 測量不精確\newline 血管張力變化時不可靠 \\
\midrule
Internally calibrated PWA & CO 測量精確\newline 連續 & 急性血管張力變化時不可靠 \\
\midrule
Externally calibrated PWA & 可靠\newline 額外血流動力學變數 & 血管張力變化後需重新校正\newline 需專用動脈導管 \\
\midrule
Right-sided thermodilution & 可靠\newline 完整血流動力學評估\newline RV 功能評估 & 侵入性\newline 可能為間歇性\newline SVV/PLR 不適用於 fluid responsiveness 預測 \\
\bottomrule
\end{longtable}

%============================================================
\section{如何選擇監測技術}
%============================================================

\subsection{選擇原則}

\begin{enumerate}
    \item 考量\textbf{侵入性}
    \item 考量\textbf{患者狀況}及額外變數的價值
    \item Echocardiography 可提供完整心臟評估，但難以每日重複超過兩次，不適合作為複雜病例的唯一監測工具
\end{enumerate}

\subsection{適應症建議}

\begin{longtable}{p{5cm}p{9cm}}
\toprule
\textbf{臨床情境} & \textbf{建議監測工具} \\
\midrule
\endfirsthead
\toprule
\textbf{臨床情境} & \textbf{建議監測工具} \\
\midrule
\endhead
Right heart / RV dysfunction & \textcolor{emergency_blue}{Pulmonary artery catheter}\newline 可測量 pulmonary artery pressures \\
\midrule
Cardiogenic shock + VA-ECMO & \textcolor{emergency_blue}{Pulmonary artery catheter}\newline 評估左心腔 loading conditions\newline 連續 SvO\textsubscript{2} 監測 \\
\midrule
Sepsis（尤其合併 ARDS） & \textcolor{emergency_blue}{Transpulmonary thermodilution}\newline 可評估 lung water and permeability \\
\midrule
圍術期（預期顯著 preload 變化） & \textcolor{emergency_blue}{Non-calibrated pulse wave analysis}\newline 即使非侵入性也可使用 \\
\bottomrule
\end{longtable}

%============================================================
\section{如何解讀 CO 測量}
%============================================================

\subsection{CO 測量的臨床價值}

\begin{enumerate}
    \item \textbf{診斷 shock 類型}：Distributive vs. others
    \item \textbf{觸發治療}：需整合其他變數評估 CO 適當性
    \begin{itemize}
        \item 臨床評估
        \item SvO\textsubscript{2}
        \item Veno-arterial PCO\textsubscript{2} difference
        \item Capillary refill time
        \item Lactate
    \end{itemize}
    \item \textbf{評估治療效果}：監測血流動力學介入或自發生理變化的影響
\end{enumerate}

\note{
在 septic shock 中，microcirculation 障礙可能使 CO 變化與組織氧合變化「脫鉤」(dissociate)，需特別注意。
}

\subsection{Least Significant Change}

CO 變化必須超過測量誤差才具臨床意義。此閾值稱為 \textbf{least significant change}，定義為可靠區分於誤差的最小變化量。

\begin{longtable}{p{5cm}p{9cm}}
\toprule
\textbf{測量方法} & \textbf{Least Significant Change} \\
\midrule
\endfirsthead
\toprule
\textbf{測量方法} & \textbf{Least Significant Change} \\
\midrule
\endhead
Pulse-derived CO & 1-2\% \\
\bottomrule
\end{longtable}

Pulse-derived CO 的高精確度使其適合評估小幅 CO 變化，如 preload responsiveness tests（mini-fluid challenge、end-expiratory occlusion test）。

%============================================================
\section{Key Points}
%============================================================

\begin{enumerate}
    \item CO 監測適用於初始復甦後仍持續循環衰竭的患者
    \item Echocardiography 可提供完整心臟評估，但為間歇性，不適合作為複雜病例的唯一監測
    \item Non-calibrated pulse wave analysis 在血管張力急性變化時不可靠
    \item Transpulmonary thermodilution 校正的裝置更準確，並提供額外血流動力學變數
    \item PAC 適用於 RV dysfunction 或 cardiogenic shock with VA-ECMO
    \item Transpulmonary thermodilution 適用於 sepsis 尤其合併 ARDS
    \item CO 變化需超過 least significant change（pulse-derived CO 為 1-2\%）才具臨床意義
    \item CO 適當性需整合臨床評估、SvO\textsubscript{2}、lactate 等多項指標判斷
\end{enumerate}

%============================================================
\section{參考文獻}
%============================================================

\begin{enumerate}
    \item De Backer D, Hajjar L, Monnet X. Monitoring cardiac output. \href{https://doi.org/10.1007/s00134-024-07566-6}{\textit{Intensive Care Med}. 2024;50:1912-1915.}

    \item De Backer D, Vieillard-Baron A. Clinical examination: a trigger but not a substitute for hemodynamic evaluation. \href{https://doi.org/10.1007/s00134-018-5489-y}{\textit{Intensive Care Med}. 2019;45(2):269-271.}

    \item Teboul JL, Saugel B, Cecconi M, et al. Less invasive hemodynamic monitoring in critically ill patients. \href{https://doi.org/10.1007/s00134-016-4375-7}{\textit{Intensive Care Med}. 2016;42(9):1350-1359.}

    \item Vincent JL, De Backer D. Circulatory shock. \href{https://doi.org/10.1056/NEJMra1208943}{\textit{N Engl J Med}. 2013;369(18):1726-1734.}

    \item Kouz K, Scheeren TWL, de Backer D, Saugel B. Pulse wave analysis to estimate cardiac output. \href{https://doi.org/10.1097/ALN.0000000000003553}{\textit{Anesthesiology}. 2021;134:119-126.}
\end{enumerate}

\end{document}
